% Workarounds for lwarp problems, used by the Checker Framework manual and by
% the Annotation File Utilities manual.  Input this file after loading lwarp.
% It affects only the HTML version.

\begin{warpHTML}

% lwarp escapes < within Verbatim, but not in ordinary text, \texttt, or
% alltt.  A browser would treat text such as "List<Object>" as an HTML tag and
% not display it.  So, within the document, < is an active character that
% produces an HTML entity.  (An unescaped > is legal in HTML text.  Leaving >
% alone lets it continue to delimit arguments, as in the manual's \<...>.)
\begingroup
\catcode`\<=\active
\gdef<{\ifmmode\string<\else\textless\fi}
\endgroup
\AfterEndPreamble{\catcode`\<=\active}
% \@noligs, which alltt calls, redefines < to print itself.
% Ligatures are irrelevant in HTML.
\makeatletter
\def\verbatim@nolig@list{\do\`\do\>\do\,\do\'\do\-}
\makeatother

% BibTeX's alpha style defines \etalchar using math mode, which lwarp would
% convert to an image, and which fails within a citation.
\newcommand{\etalchar}[1]{\textsuperscript{#1}}
% The .bbl file uses \newcommand to define \etalchar.
\pretocmd{\bibliography}{\let\etalchar\relax}{}{}
\AtBeginEnvironment{thebibliography}{\renewcommand{\etalchar}[1]{\textsuperscript{#1}}}

\end{warpHTML}
